\section{Documento descriptivo y especificaciones de funcionamiento} \label{sec: Sec1}

\subsection{Objetivos generales del proyecto} \label{subsec: SubSec11}

El proyecto tiene como propósito la obtención de un \textbf{líquido de salida, S}. El proceso comenzará con la inyección de un \textbf{líquido de entrada, E} en un tanque. Este líquido, de viscosidad cercana a la del agua, llegará a través de una válvula, y será precalentado en una zona de la tubería de alimentación antes de disponerse en el tanque.

Una vez el tanque esté en su nivel máximo, se dejará de introducir \textbf{E}, y comenzará el proceso de calentamiento y agitación a baja velocidad. 

Durante este proceso, llegado a un cierto nivel de temperatura, el líquido comenzará a reaccionar, pasándo a estado gaseoso en forma de \textbf{vapor, V}, que deberá ser expulsado para no superar un umbral de presión.

Una vez el tanque esté a mitad de capacidad, se dejará de aportar calor mientras se sigue agitando, hasta llegar a una temperatura umbral, donde se prescindirá del agitador.

La extracción del tanque se llevará a cabo de forma cautelar y bajando su temperatura, manteniendo unos niveles de presión y temperatura adecuados.

El líquido se va llevando a un segundo reactor controlando la concentración del mismo, a la vez que se aporta otro elemento. En este segundo reactor se mantendrá una temperatura constante.

Terminada esta reacción, se vacía el segundo reactor y el producto final \textbf{S} se almacena.
